\documentclass{article}
\usepackage[T1]{fontenc}
\usepackage{lmodern}
\usepackage{graphicx}
\usepackage{amsmath,amssymb}
\usepackage[a4paper,margin=2cm]{geometry}
\usepackage[absolute,overlay]{textpos}
\setlength{\TPHorizModule}{1mm}
\setlength{\TPVertModule}{1mm}
\usepackage[indonesian]{babel}
\usepackage{hyperref}
\setlength{\emergencystretch}{2em}
% unit: Halaman 12

\begin{document}

% === Halaman 12 (python/native) ===

Contoh: Kurva eliptik y2 = x3 + 2x + 4 

   Misalkan P(2, 4) dan Q(0, 2) dua titik pada kurva

   Penjumlahan titik: P + Q = R. Tentukan R!

   Langkah-langkah menghitung koordinat R:

• Gradien garis g: m = (yp – yq)/(xp – xq) =(4 – 2)/(2 – 0) = 1

• xr = m2 – xp – xq  = 12 – 2  – 0 = –1   

• yr = m(xp – xr) – yp = 1(2 – (-1)) – 4 = –1

• Jadi koordinat R(-1, -1) 

• Periksa apakah R(-1, -1) sebuah titik pada kurva eliptik:

     y2 = x3 + 2x + 4    (-1)2 = (-1)3 + 2(-1) + 4



    1 = -1 – 2 + 4 



    1 = 1   (terbukti R(-1,-1) titik yang terletak pada 




         kurva y2 = x3 + 2x + 4 ) 


12

\end{document}
